\documentclass[10pt]{beamer}
\usetheme{metropolis}

\usepackage{amsmath,amssymb}
\usepackage{mathtools}
\usepackage{xeCJK}
\IfFontExistsTF{Noto Serif SC}{\setCJKmainfont{Noto Serif SC}}{\setCJKmainfont{Noto Serif CJK SC}}
\usepackage{tikz-cd}

% 1. 开启标题背景色填充
\metroset{block=fill}

% theorem environments
\newtheorem{proposition}{Proposition}
\newtheorem{exercise}{Exercise}

% notation
\newcommand{\R}{\mathbb{R}}
\newcommand{\N}{\mathbb{N}}
\newcommand{\F}{\mathcal{F}}
\newcommand{\supp}{\operatorname{supp}}
\newcommand{\vX}{\mathfrak{X}}
\newcommand{\D}{\mathcal{D}}


\title{9\quad Vector Fields}
\subtitle{Smoothness, Lie Brackets \& Integral Curves}
\date{October 29, 2026}
\author{LXJ}

\begin{document}

\maketitle

%====================================================================
\section{Vector Fields}

\begin{frame}{Smooth curves on a closed interval}
	\begin{definition}[Smooth curve]
		A mapping $\sigma : [a, b] \to M$ is a \emph{smooth curve} in
		$M$ if $\sigma$ extends to a $C^\infty$-map of
		$(a - \varepsilon, b + \varepsilon) \to M$ for some
		$\varepsilon > 0$.
	\end{definition}
	\begin{block}{Note}
		If $\sigma : [a, b] \to M$ is a smooth curve, then its tangent
		vector
		\[
			\dot{\sigma}(t)
			= \sigma_*\Bigl( \frac{d}{dt} \Big|_{t} \Bigr)
			\in T_{\sigma(t)} M
		\]
		is well-defined for all $t \in [a, b]$.
	\end{block}
\end{frame}

\begin{frame}[fragile]{Vector fields along a curve}
	\begin{definition}[Vector field along a curve]
		A smooth vector field $X$ along a curve
		$\sigma : [a, b] \to M$ is a $C^\infty$-map
		\[
			X : [a, b] \longrightarrow TM
		\]
		which \emph{lifts} $\sigma$, i.e.
		\[
			\pi \circ X = \sigma .
		\]
	\end{definition}
	\medskip
	\begin{center}
		\begin{tikzcd}
			& TM \arrow[d, "\pi"] \\
			{[a,b]} \arrow[ru, "X"] \arrow[r, "\sigma"'] & M
		\end{tikzcd}
	\end{center}
\end{frame}

\begin{frame}{Vector fields on an open subset}
	\begin{definition}[Vector fields on an open subset of $M$]
		A smooth vector field on an open set $U$ in $M$ is a
		$C^\infty$-map $X : U \to TM$ s.t.
		\[
			\pi \circ X = \mathrm{id}_U .
		\]
	\end{definition}
	Sometimes we call a smooth vector field $X$ on $M$ a
	\emph{section} of $TM$. The set of all smooth sections on $U$ of
	$TM$ is denoted by $\Gamma(U; TM)$. It is a vector space over
	$\R$ and also a $C^\infty(U)$-\emph{module}:
	\[
		\Gamma(U; TM) = C^\infty(U, TM) = \vX(U).
	\]
\end{frame}

%====================================================================
\section{A Criterion for Smoothness}

\begin{frame}{Three equivalent conditions}
	\begin{proposition}
		Let $X$ be a vector field on $M$. Then the following are
		equivalent:
		\begin{enumerate}
			\item[(a)] $X$ is $C^\infty$;
			\item[(b)] if $(U, x_1, \dots, x_d)$ is a local coordinate
			      system on $M$ and $\{a_i\}$ are the functions on $U$
			      defined by
			      \[
				      X\big|_U = \sum_{i=1}^{d} a_i\,
				      \frac{\partial}{\partial x_i},
			      \]
			      then $a_i \in C^\infty(U)$;
			\item[(c)] whenever $V \subset M$ is open, if
			      $f \in C^\infty(V)$, then $X(f) \in C^\infty(V)$.
		\end{enumerate}
	\end{proposition}
\end{frame}

\begin{frame}{Proof: (a) $\Rightarrow$ (b)}
	Let $X : M \to TM$ be a $C^\infty$-map s.t.\ $\pi \circ X =
		\mathrm{id}_M$, so $X(p) = V_p \in T_p M$. Take a local coordinate
	system $(U, \varphi, x_1, \dots, x_d)$ about $p \in M$; the local
	coordinate system about $V_p$ is given by the local trivialization
	$\widetilde{\varphi} : \pi^{-1}(U) \to \R^{2d}$.
	\medskip

	Smoothness of $X$ in local coordinates requires
	$\widetilde{\varphi} \circ X \circ \varphi^{-1}$ to be $C^\infty$
	from $\varphi(U) \subset \R^d$ to
	$\widetilde{\varphi}\bigl(\pi^{-1}(U)\bigr) \subset \R^{2d}$.
	Writing $X(p) = V_p = \sum_i a_i(p) \left. \frac{\partial}{\partial
			x_i} \right|_{p}$,
	\begin{align*}
		\widetilde{\varphi} \circ X \circ \varphi^{-1}(x_1, \dots, x_d)
		 & = \Bigl( \varphi\bigl(\varphi^{-1}(x_1, \dots, x_d)\bigr);\
		a_i\bigl(\varphi^{-1}(x_1, \dots, x_d)\bigr) \Bigr)                      \\
		 & = \bigl( x_1, \dots, x_d;\ a_i(\varphi^{-1}(x_1, \dots, x_d)) \bigr),
	\end{align*}
	which means $a_i \in C^\infty(U)$.
\end{frame}

\begin{frame}{Proof: (b) $\Rightarrow$ (c) $\Rightarrow$ (a)}
	\textbf{(b) $\Rightarrow$ (c).} For $f \in C^\infty(V)$ we show
	$X(f)$ is $C^\infty$ about each $p \in V$. Take a local coordinate
	system $(U, x_1, \dots, x_d)$ about $p$; then
	$X\big|_U = \sum_i a_i\, \frac{\partial}{\partial x_i}$ with
	$a_i \in C^\infty(U)$, so
	\[
		X(f) = \sum_{i=1}^{d} a_i\, \frac{\partial f}{\partial x_i}
		= \sum_{i=1}^{d} a_i\,
		\frac{\partial (f \circ \varphi^{-1})}{\partial x_i}
		\in C^\infty(U).
	\]
	\textbf{(c) $\Rightarrow$ (a).} In an arbitrary local coordinate
	system $(U, x_1, \dots, x_d)$, write
	$X = \sum_i a_i\, \frac{\partial}{\partial x_i}$. Apply (c) to the
	coordinate function $x_i \in C^\infty(U)$:
	\[
		X(x_i) = a_i \in C^\infty(U). \qedhere
	\]
\end{frame}

%====================================================================
\section{The Lie Bracket}

\begin{frame}{The Lie bracket of two vector fields}
	If $X, Y$ are smooth vector fields on $M$ (i.e.\ $X, Y \in
		\vX(M) = \Gamma(TM)$), we define a vector field $[X, Y]$,
	the \emph{Lie bracket} of $X$ and $Y$, by
	\[
		[X, Y]_p(f) = X_p\bigl(Y(f)\bigr) - Y_p\bigl(X(f)\bigr).
	\]
	\medskip
	\begin{block}{Note}
		The bracket measures the \emph{failure} of the flows of $X$ and
		$Y$ to commute; locally it is the commutator of the derivations
		$X$ and $Y$.
	\end{block}
\end{frame}

\begin{frame}{Properties of the Lie bracket}
	\begin{proposition}
		Let $X, Y, Z \in \vX(M)$ and $f, g \in C^\infty(M)$.
		\begin{enumerate}
			\item[(a)] $[X, Y]$ is a smooth vector field on $M$;
			\item[(b)] $[fX, gY] = fg\,[X, Y] + f\,(Xg)\, Y
				      - g\,(Yf)\, X$;
			\item[(c)] $[X, Y] = -[Y, X]$
			      \emph{(skew-symmetry)};
			\item[(d)] $\bigl[[X, Y], Z\bigr] + \bigl[[Y, Z], X\bigr]
				      + \bigl[[Z, X], Y\bigr] = 0$
			      \emph{(Jacobi identity)}.
		\end{enumerate}
	\end{proposition}
	\begin{block}{Note (Lie algebra)}
		A vector space with a bilinear operation satisfying (c) and (d)
		is called a \emph{Lie algebra}.
	\end{block}
\end{frame}

%====================================================================
\section{Integral Curves}

\begin{frame}{Integral curves}
	\begin{definition}[Integral curve]
		Let $X \in \vX(M)$. A smooth curve $\sigma$ in $M$ is an
		\emph{integral curve} of $X$ if
		\[
			\dot{\sigma}(t) = X\bigl(\sigma(t)\bigr)
			\qquad (\forall\, t \in I).
		\]
	\end{definition}
	\medskip
	\begin{block}{Reformulation}
		A curve $\gamma : (a, b) \to M$ is an integral curve of $X$
		$\iff$
		\[
			\gamma_*\Bigl( \frac{d}{dt} \Big|_{t} \Bigr)
			= X\bigl(\gamma(t)\bigr)
			\qquad (t \in (a, b)).
		\]
	\end{block}
\end{frame}

\begin{frame}{Existence and uniqueness: local form}
	Suppose $0 \in (a,b)$ and $\gamma(0) = p$. Choose a local coordinate
	system $(U, \varphi)$ with coordinate functions $x_1, \dots, x_d$
	about $p$; write
	\[
		X\big|_U = \sum_{i=1}^{d} f_i\, \frac{\partial}{\partial x_i}
		\qquad (f_i \in C^\infty(U)).
	\]
	For all $t$ with $\gamma(t) \in U$,
	\[
		\gamma_*\Bigl( \frac{d}{dt} \Big|_{t} \Bigr)
		= \sum_{i=1}^{d} \frac{d (x_i \circ \gamma)}{dt} \Big|_{t}
		\left. \frac{\partial}{\partial x_i} \right|_{\gamma(t)}
		= X\bigl(\gamma(t)\bigr)
		= \sum_{i=1}^{d} f_i\bigl(\gamma(t)\bigr)
		\left. \frac{\partial}{\partial x_i} \right|_{\gamma(t)} .
	\]
\end{frame}

\begin{frame}{Reduction to a system of ODEs}
	Thus $\gamma$ is an integral curve of $X$ on $\gamma^{-1}(U)$ iff,
	with $\gamma_i \triangleq x_i \circ \gamma$,
	\[
		(\star)\qquad
		\frac{d \gamma_i}{dt} \Big|_{t}
		= f_i\bigl( \varphi^{-1}\bigl( \gamma_1(t), \dots,
		\gamma_d(t) \bigr) \bigr)
		\qquad (t \in \gamma^{-1}(U)).
	\]
	\medskip
	\begin{alertblock}{Theorem (Existence \& uniqueness)}
		$(\star)$ is a system of ODEs, for which there exists a
		fundamental existence \& uniqueness theorem; hence integral
		curves of a smooth vector field exist and are unique.
	\end{alertblock}
\end{frame}

\section{Existence and Uniqueness}

\begin{frame}{Integral curves: the question}
	\begin{definition}[Integral curve]
		Let $X$ be a smooth vector field on $M$. A smooth curve
		$\sigma$ in $M$ is an \emph{integral curve} of $X$ if
		\[
			\dot{\sigma}(t) = X\bigl(\sigma(t)\bigr)
		\]
		for each $t$ in the domain of $\sigma$.
	\end{definition}
	\medskip
	\begin{block}{Question}
		Given a smooth vector field on $M$ and a point $p \in M$, does
		there exist an integral curve of $X$ through $p$? If so, is
		there a unique one?
	\end{block}
\end{frame}

\begin{frame}{The ODE system in local coordinates}
	We interpret in coordinates! Let $0 \in (a,b)$, $p = \sigma(0)$.
	Let $(U, \varphi = (x_1, \dots, x_d))$ be a local coordinate system
	(centered at $p$). Then the vector field $X$ can be written as
	\[
		X = \sum_{i=1}^{d} f_i\, \frac{\partial}{\partial x_i},
		\qquad f_i \in C^\infty(U).
	\]
	Assuming $\sigma\bigl((a,b)\bigr) \subset U$,
	\[
		\dot{\sigma}(t) = \sigma_*\Bigl( \frac{d}{dt} \Big|_{t} \Bigr)
		= \sum_{i=1}^{d} \frac{d (x_i \circ \sigma)}{dt}
		\left. \frac{\partial}{\partial x_i} \right|_{\sigma(t)}
		\quad \forall\, t \in (a,b).
	\]
	So $\dot{\sigma}(t) = X\bigl(\sigma(t)\bigr)$ iff
	\[
		\frac{d \sigma_i}{dt}
		= f_i\bigl( \varphi^{-1}\bigl( \sigma_1(t), \dots,
		\sigma_d(t) \bigr) \bigr),
		\quad i = 1, \dots, d . \tag{1}
	\]
\end{frame}

\begin{frame}{Reduction to ODEs}
	Here $\bigl( \sigma_1, \dots, \sigma_d \bigr)
		= \bigl( x_1 \circ \sigma, \dots, x_d \circ \sigma \bigr)
		\in \varphi(U)$ are the coordinates of the point
	$\sigma(t) \in M$.
	\medskip

	\emph{(1) is a system of ODEs, for which there exist fundamental
		existence and uniqueness theorems. In terms of the terminology for
		manifolds, we have:}
	\medskip

	\begin{theorem}
		Let $X$ be a $C^\infty$ vector field on $M$. For each $p \in M$
		there exist $a(p)$, $b(p) \in \R \cup \{\pm\infty\}$ and a
		smooth curve $\sigma_p : \bigl( a(p), b(p) \bigr) \to M$ such
		that
		\begin{enumerate}
			\item[(a)] $0 \in \bigl( a(p), b(p) \bigr)$ and
			      $\sigma_p(0) = p$;
			\item[(b)] $\sigma_p$ is an integral curve of $X$;
			\item[(c)] if $\mu : (c, d) \to M$ satisfies (a), (b), then
			      $(c, d) \subset \bigl( a(p), b(p) \bigr)$ and
			      $\sigma_p\big|_{(c,d)} = \mu$.
		\end{enumerate}
	\end{theorem}
\end{frame}

\begin{frame}{Proof of uniqueness and maximality}
	If $\alpha, \beta$ are integral curves defined on open intervals
	$I$, $J$, assume they have the same initial condition
	$\alpha(t_0) = \beta(t_0)$ for the same $t_0 \in I \cap J$. Then
	$I \cap J \neq \varnothing$. Define
	\[
		K = \bigl\{ t \in I \cap J \bigm| \alpha(t) = \beta(t) \bigr\}.
	\]
	\begin{itemize}
		\item $K$ is \emph{open} (by the uniqueness theorem of ODE);
		\item $K$ is \emph{closed} (by continuity);
		\item hence $K = I \cap J$ (by the connectedness of
		      $I \cap J$).
	\end{itemize}
	\medskip
	Thus (c) is true --- (a), (b) as well: $\bigl( a(p), b(p) \bigr)$
	is the union of all domains containing $t = 0$ on which an integral
	curve with $\sigma(0) = p$ exists, and it is non-empty by the
	fundamental existence theorem.
\end{frame}

\section{The Flow of a Vector Field}

\begin{frame}{The transformations $X_t$}
	\begin{definition}[$X_t$ and its domain]
		$\forall\, t \in \R$, we define a transformation $X_t$ with
		domain
		\[
			\D_t = \bigl\{ p \in M \bigm|
			t \in \bigl( a(p), b(p) \bigr) \bigr\},
		\]
		by setting
		\[
			X_t(p) = \sigma_p(t).
		\]
	\end{definition}
	\medskip
	\begin{center}
		\emph{$X_t$ moves each point $p$ for time $t$ along its
			maximal integral curve.}
	\end{center}
\end{frame}

\begin{frame}{Properties of the flow}
	\begin{theorem}[Properties of the flow]
		\begin{enumerate}
			\item[(d)] $\forall\, p \in M$, $\exists$ a neighborhood $V$
			      of $p$ and $\varepsilon > 0$ s.t.\ the map
			      $(t, p) \longmapsto X_t(p)$ is defined and $C^\infty$
			      from $(-\varepsilon, \varepsilon) \times V \to M$;
			\item[(e)] $\D_t$ is open for each $t$;
			\item[(f)] $\bigcup_{t > 0} \D_t = M$;
			\item[(g)] $X_t : \D_t \longrightarrow \D_{-t}$ is a
			      diffeomorphism with inverse $X_{-t}$;
			\item[(h)] let $s, t \in \R$. Then the domain of $X_s \circ X_t$
			      is contained in but generally not equal to $\D_{s+t}$.
			      However, the domain of $X_s \circ X_t$ \emph{is}
			      $\D_{s+t}$ in the case in which $s, t$ have the same
			      sign. And we have
			      \[
				      X_s \circ X_t = X_{s+t}.
			      \]
		\end{enumerate}
	\end{theorem}
\end{frame}

%====================================================================
\section{Proof of the Flow Properties}

\begin{frame}{Smooth dependence and (f)}
	\textbf{(d)} holds true according to the ODE theorem that the
	solution depends on both $t$ and the initial value $\varphi(p)$
	smoothly. So there's a neighborhood $V$ of $p$ and
	$(-\varepsilon, \varepsilon)$ of $t = 0$ s.t.\ $\D_t$ is defined
	for all $t \in (-\varepsilon, \varepsilon)$ and $V \subset \D_t$.
	\bigskip

	\textbf{(f)} Let $\bigl( a(p), b(p) \bigr)$ be the union of the
	domains of $t$ that contain $t = 0$ and on which $\sigma(t)$ is an
	integral curve of $X$ with $\sigma(0) = p$. By the fundamental
	existence theorem, $\bigl( a(p), b(p) \bigr) \neq \varnothing$.
	Therefore
	\[
		\bigcup_{t > 0} \D_t = M . \qed
	\]
\end{frame}

\begin{frame}{Proof of (h): shifting the initial point}
	Let $t \in \bigl( a(p), b(p) \bigr)$. Then
	$s \longmapsto \sigma_p(t + s)$ is an integral curve of $X$
	initiated at $\sigma_p(t)$, with maximal domain
	$\bigl( a(p) - t,\, b(p) - t \bigr)$. By (c),
	\[
		\bigl( a(p) - t,\, b(p) - t \bigr)
		= \bigl( a\bigl(\sigma_p(t)\bigr),\,
		b\bigl(\sigma_p(t)\bigr) \bigr).
	\]
	So
	\[
		\sigma_p(t + s) = \sigma_{\sigma_p(t)}(s). \tag{2}
	\]
	Now, if $p \in$ the domain of $X_s \circ X_t$, then
	\[
		t \in \bigl( a(p), b(p) \bigr)
		\quad \text{and} \quad
		s \in \bigl( a\bigl(\sigma_p(t)\bigr),\,
		b\bigl(\sigma_p(t)\bigr) \bigr)
		= \bigl( a(p) - t,\, b(p) - t \bigr).
	\]
	So $s + t \in \bigl( a(p), b(p) \bigr)$, i.e.\ $p \in \D_{s+t}$,
	and we know $X_s \circ X_t = X_{s+t}$ by (2).
\end{frame}

\begin{frame}{Opposite signs: the domain can shrink}
	For $s$ and $t$ that do not have the same sign, the inclusion may
	be strict. \textbf{Eg.} On $\R^2 \setminus \{0\}$, consider
	$\frac{\partial}{\partial r_1}$ and take $s = -1$, $t = 1$. Then
	the domain of
	\[
		X_{-1} \circ X_1 \quad \text{is} \quad
		\R \setminus [-1, 0],
	\]
	but
	\[
		\D_{-1+1} = \D_0 = \R \setminus \{0\}.
	\]
	\medskip
	\begin{block}{Note}
		The composed map can be defined \emph{less} often than
		$X_{s+t}$: a point may leave the manifold at time $t$ and the
		way back at time $s$ need not exist.
	\end{block}
\end{frame}

\begin{frame}{Same sign: the domain of $X_s \circ X_t$ is $\D_{s+t}$}
	If, however, $s, t$ have the same sign, for eg.\ $s, t > 0$:
	$\forall\, p \in \D_{s+t}$,
	\[
		s + t \in \bigl( a(p), b(p) \bigr)
		\ \Rightarrow\
		\left.\begin{aligned}
			t & \in \bigl( a(p), b(p) \bigr)           \\
			s & \in \bigl( a(p) - t,\, b(p) - t \bigr)
			= \bigl( a\bigl(\sigma_p(t)\bigr),\,
			b\bigl(\sigma_p(t)\bigr) \bigr)
		\end{aligned}\right\}
	\]
	\[
		\Rightarrow\
		p \in \D_t,\quad \sigma_p(t) \in \D_s
		\ \Rightarrow\
		p \in \text{domain of } X_s \circ X_t .
	\]
	\medskip
	\textbf{(e)} is trivial when $t = 0$, i.e.\ $\D_0 = M$ is open.
	Assume that $t > 0$ and $p \in \D_t$\,\dots
\end{frame}

%====================================================================
\section{Completion of the Flow Proofs}

\begin{frame}{Recall (h): shifting the initial point}
	Let $t, s \in \R$. For any $p$ in the domain of $X_s \circ X_t$ we
	see that
	\[
		p \in \D_t
		\quad \bigl( t \in \bigl( a(p), b(p) \bigr) \bigr)
		\quad \text{and} \quad
		\sigma_p(t) = X_t(p) \in \D_s .
	\]
	To see that $p \in \D_{s+t}$, notice that we can define a
	$C^\infty$-map $\gamma : s \longmapsto \sigma_p(t + s)$. Then
	\[
		\gamma(0) = \sigma_p(t), \qquad
		\dot{\gamma}(s) = X\bigl(\sigma_p(t+s)\bigr) = X\bigl(\gamma(s)\bigr).
	\]
	So $\bigl( a(\sigma_p(t)), b(\sigma_p(t)) \bigr)$ is the maximal
	interval for $\gamma$, hence equals $\bigl( a(p) - t,\, b(p) - t
	\bigr)$, and
	\[
		\sigma_p(t + s) = \sigma_{\sigma_p(t)}(s). \tag{2}
	\]
	So if $p \in \D_t$ and $\sigma_p(t) \in \D_s$ then
	$p \in \D_{t+s}$. The inverse \emph{doesn't} hold! For eg.\ the
	domain of $X_{-0.9} \circ X_1$ is contained in $\D_{0.1}$ but not
	equal to it.
\end{frame}

\begin{frame}{Proof of (e): $\D_t$ is open --- setup}
	Obviously true when $t = 0$, since $\D_0 = M$. Fix $t > 0$
	(similar for $t < 0$) and $p \in \D_t$. We need a neighborhood
	$N_p$ s.t.\ $N_p \subset \D_t$, i.e.\ $\forall\, q \in N_p$,
	$t \in \bigl( a(q), b(q) \bigr)$.
	\medskip

	Consider the curve $\sigma_p\bigl( [0, t] \bigr)$, which is
	\emph{compact} in $M$, and let
	\[
		\Gamma = \bigl\{ \bigl( \tau, \sigma_p(\tau) \bigr)
		\bigm| \tau \in [0, t] \bigr\}.
	\]
	By (d), for each point $\bigl( t_0, \sigma_p(t_0) \bigr) \in
	\Gamma$ there is an $\varepsilon_{t_0} > 0$ and
	$V_{\sigma_p(t_0)}$ open about $\sigma_p(t_0)$ s.t.\
	$(-\varepsilon_{t_0}, \varepsilon_{t_0}) \times
	V_{\sigma_p(t_0)}$ carries the flow.
\end{frame}

\begin{frame}{Proof of (e): compactness and induction}
	By the compactness of $\Gamma$, there is a finite subcover for
	$\sigma_p\bigl( [0, t] \bigr)$: an interval $(-\varepsilon,
	\varepsilon)$ and a neighborhood $W$ of $\sigma_p\bigl( [0,t]
	\bigr)$ s.t.\ the flow is defined on $(-\varepsilon, \varepsilon)
	\times W$. Let $n$ be an integer $> 1$ s.t.\ $\bigl| \frac{t}{n}
	\bigr| < \varepsilon$. Define
	\[
		W_1 = \bigl\{ p \in W \bigm| X_{\frac{t}{n}}(p) \in W \bigr\}
		= W \cap X_{\frac{t}{n}}^{-1}(W),
	\]
	which is open in $W$; inductively,
	\[
		W_k = W \cap X_{\frac{t}{n}}^{-1}(W_{k-1}) \subset W
		\quad \text{open}, \qquad p \in W_n .
	\]
	$W_n$ consists of all points remaining in $W$ after $n$
	translations, so $W_n \subset \D_t$. Therefore $\D_t$ is open.
\end{frame}

\begin{frame}{Proof of (g): $X_t$ is a diffeomorphism}
	$X_t : \D_t \longrightarrow \D_{-t}$ with inverse $X_{-t}$. It
	follows from
	\[
		X_t = X_{\frac{t}{n}} \circ \cdots \circ X_{\frac{t}{n}},
	\]
	with each $X_{\frac{t}{n}}$ $C^\infty$ on a neighborhood of
	$p \in \D_t$, that $X_t$ is $C^\infty$.
	\medskip

	$X_t$ is $1{:}1$ because
	\[
		X_{-t} \circ X_t(p) = \sigma_{\sigma_p(t)}(-t) = p
	\]
	by the uniqueness of the initial-value problem of ODE. And
	$X_{-t}$ is also $C^\infty$.
\end{frame}

%====================================================================
\section{Complete Vector Fields}

\begin{frame}{Complete vector fields and 1-parameter groups}
	\begin{definition}[Complete vector field]
		A smooth vector field $X$ on $M$ is \emph{complete} if
		$\D_t = M$ for all $t$ (i.e.\ all the points of $M$ can be
		translated up till time $t$).
	\end{definition}
	\begin{itemize}
		\item The set of transformations $\{ X_t \mid t \in \R \}$
		      forms a \emph{group} of transformations, isomorphic to
		      $(\R, +)$, because $X_t \circ X_s = X_{t+s}$.
		\item If $X$ is not complete, the transformations $X_t$ do
		      \emph{not} form a group since the domain may depend on
		      $t$. In this case we refer to the collection
		      $\{ X_t \}$ as the \emph{local 1-parameter group}
		      of $X$.
	\end{itemize}
\end{frame}

\begin{frame}{Remark: a simple noncomplete example}
	Consider $\frac{\partial}{\partial r_1}$ on $\R^2 \setminus \{0\}$.
	The integral curve $\sigma$ passing through $(a, b)$ is given by
	\[
		\left\{
		\begin{aligned}
			\dot{\sigma}(t)
			 & = \left. \frac{\partial}{\partial r_1}
			\right|_{\sigma(t)} \\
			\sigma(0)          & = (a, b)
		\end{aligned}
		\right.
		\iff
		\left\{
		\begin{aligned}
			\frac{d r_1}{dt} & = 1 \\
			\frac{d r_2}{dt} & = 0
		\end{aligned}
		\right.
		\quad (r_1, r_2)(0) = (a, b).
	\]
	So the integral curve starting from $(a, 0)$ ($a > 0$) has a life
	span $(-a, +\infty)$.
	\medskip

	\begin{block}{Note}
		If the manifold $M$ is \emph{compact}, any $C^\infty$ vector
		field on $M$ is complete! (Exercise !!)
	\end{block}
\end{frame}

%====================================================================
\section{Straightening a Vector Field}

\begin{frame}{Prop 1.53: straightening out a nonzero vector field}
	\begin{alertblock}{Proposition 1.53}
		Let $p \in M^d$, and let $X$ be a smooth vector field on $M$
		s.t.\ $X(p) \neq 0$. Then there is a coordinate system
		$(U, \varphi) = (U, x_1, \dots, x_d)$ about $p$ s.t.
		\[
			X\big|_U = \left. \frac{\partial}{\partial x_1} \right|_U .
		\]
	\end{alertblock}
	\medskip
	\begin{center}
		\emph{Near a non-vanishing point, $X$ can be taken as the
			first coordinate vector field.}
	\end{center}
\end{frame}

\begin{frame}{Proof: the flow map}
	Choose a coordinate system centered at $p$, denoted by
	$(V, \psi) = (V, y_1, \dots, y_d)$, s.t.
	\[
		X\big|_p = \left. \frac{\partial}{\partial y_1} \right|_{p}.
	\]
	Consider the flow defined by
	\[
		\left\{
		\begin{aligned}
			\dot{\sigma}(t) & = X\bigl(\sigma(t)\bigr), \\
			\sigma(0)       & = \psi^{-1}(0, \eta_2, \dots, \eta_d).
		\end{aligned}
		\right.
	\]
	Then $\exists\, \varepsilon > 0$ and a neighborhood
	$W \subset \R^{d-1}$ about $0$ s.t.
	\[
		\sigma(t, \eta_2, \dots, \eta_d)
		= X_t\bigl( \psi^{-1}(0, \eta_2, \dots, \eta_d) \bigr)
	\]
	is defined and smooth for
	$(t, \eta_2, \dots, \eta_d) \in (-\varepsilon, \varepsilon)
	\times W \subset \R^d$.
\end{frame}

\begin{frame}{Proof: $\sigma$ is non-singular at the origin}
	\[
		\sigma_{*,0}\Bigl( \frac{d}{dt} \Big|_{0} \Bigr)
		= X(p) = \left. \frac{\partial}{\partial y_1} \right|_{p}.
	\]
	For $k = 1, \dots, d$ and $i = 2, \dots, d$,
	\begin{align*}
		\sigma_{*,0}\Bigl( \frac{\partial}{\partial \eta_i} \Big|_{0}
		\Bigr) (y_k)
		 & = \left. \frac{d}{ds} \right|_{s=0}
		y_k \circ \sigma(0, \dots, s, \dots, 0)  \\
		 & = \left. \frac{d}{ds} \right|_{s=0}
		y_k \circ X_0\bigl( \psi^{-1}(0, \dots, s, \dots, 0) \bigr) \\
		 & = \left. \frac{d}{ds} \right|_{s=0}
		y_k \circ \psi^{-1}(0, \dots, s, \dots, 0)
		= \delta_{ki}.
	\end{align*}
	So $\sigma_{*,0}\bigl( \left. \frac{\partial}{\partial \eta_i}
		\right|_{0} \bigr) = \left. \frac{\partial}{\partial y_i}
		\right|_{p}$, hence $\sigma_{*,0}$ is non-singular at $p$.
\end{frame}

\begin{frame}{Proof: the new coordinates}
	By the inverse function theorem, $\exists\, (U, \sigma^{-1})$ about
	$p$. Denote $\sigma^{-1} = (x_1, \dots, x_d)$. Then
	$\forall\, h \in C^\infty(U)$,
	\begin{align*}
		\left. \frac{\partial}{\partial x_1} \right|_{U} (h)
		 & = \left. \frac{\partial}{\partial r_1} \right|_{0}
		(h \circ \sigma)
		= \frac{\partial}{\partial t} (h \circ \sigma)  \\
		 & = \sigma_*\Bigl( \frac{\partial}{\partial t} \Bigr) (h)
		= X(h) = X\big|_U (h). \qedhere
	\end{align*}
\end{frame}



\end{document}
